\begin{blocksection}
\question Complete the definition for \texttt{all\char`_true}, which takes in a list \texttt{lst}
and returns \texttt{True} if there are no False-y values in the list and \texttt{False} otherwise.
Make sure that your implementation is recursive. What is the name of the built-in Python function that does this?

\begin{lstlisting}
def all_true(lst):
    '''
    >>> all_true([True, 1, 'True'])
    True
    >>> all_true([1, 0, 1])
    False
    >>> all_true([])
    True
    '''
\end{lstlisting}

\begin{solution}[1in]
\begin{lstlisting}:
    if not lst:
        return True
    elif not lst[0]:
        return False
    else:
        return all_true(lst[1:])
\end{lstlisting}
\begin{lstlisting}:
    The built-in Python function all does the same thing that all_true does!
\end{lstlisting}
\end{solution}

\begin{questionmeta}
    \textbf{Teaching Tips}
    \begin{enumerate}
            \item It may be helpful to review recursion and discuss what the base case and recursive case should be.
    \end{enumerate}
\end{questionmeta}
\end{blocksection}

